\chapter{TANGENT SPACES, BUNDLES, AND VECTORS}


\section{Manifold Structure}
\noindent \textit{August 31st.}

Recall that given a smooth manifold $M$, the tangent space $T_{p}M$ at a point $p \in M$ consists of all derivations from $C_{p}^{\infty}(M)$ to $\R$. $T_{p}M$ exactly equals $T_{p}U$, where $U$ is an open neighbourhood of $p$ in $M$, since germs of functions need only match in a neighbourhood. If $M$ is a manifold of dimension $n$, then any tangent space of a point in $M$ is a vector space of dimension $n$, too. Now, the \eax{tangent bundle} of the manifold is defined as the disjoint union of all tangent spaces:
\begin{align}
    TM = \bigsqcup_{p \in M} T_{p}M.
\end{align}
From here, let $\pi:TM \to M$ denote the projection map; that is, if $v \in T_{p}M \subseteq TM$ for some $p \in M$, then $\pi(v) = p$.


We can equip $TM$ with a topology so that it becomes a smooth manifold is its own right, and so that $\pi$ also happens to be a smooth map. For $p \in M$, let $(U,\varphi)$ be a chart on $M$ containing $p$. Note that for all $p \in U$, the tangent space $T_{p}M$ has a basis $\{\frac{\partial}{\partial x^{1}}\big|_{p},\ldots,\frac{\partial}{\partial x^{n}}\big|_{p}\}$. Given any $v \in T_{p}M$, it can uniquely be written as
\begin{align}
    v = \sum_{i=1}^{n} c^{i}(v) \, \frac{\partial}{\partial x^{i}} \Big|_{p}.
\end{align}
Now, given $v \in TU$, we define a chart $(\tilde{U},\tilde{\varphi})$ as $\tilde{U} = \pi^{-1}(U) = \bigsqcup_{p \in U} T_{p}M = TU$, and
\begin{align}
    \tilde{\varphi}:TU \to \varphi(U) \times \R^{n}, \quad v \mapsto (x^{1}(p),\ldots,x^{n}(p),c^{1}(v),\ldots,c^{n}(v)) = ((x^{1} \circ \pi)(v),\ldots,(x^{n} \circ \pi)(v),c^{1}(v),\ldots,c^{n}(v)).
\end{align}
Note that $\tilde{\varphi}$ has a (set-theoretic) inverse since
\begin{align}
    \tilde{\varphi}^{-1}:\varphi(U) \times \R^{n} \to TU, \quad (\varphi(p),c^{1},\ldots,c^{n}) \mapsto v = \sum_{i=1}^{n} c^{i} \frac{\partial}{\partial x^{i}} \Big|_{p}.
\end{align}
On $\varphi(U) \times \R^{n} \subseteq \R^{2n}$, we put the subspace topology. Using the bijection $\tilde{\varphi}$, we put a topology on $TU$; that is, $A \subseteq TU$ is open if and only if $\tilde{\varphi}(A) \subseteq \varphi(U) \times \R^{n}$ is open.

With $\cU = \{(U_{\alpha},\varphi_{\alpha})\}$ is a maximal atlas of $M$, let $\cB = \{A \subseteq TU_{\alpha} \mid A \text{ open and } \alpha \text{ is any index}\}$. We show that $\cB$ is a basis for a topology on $TM$. The union being $M$ is clear since
\begin{align}
    TM = \bigcup_{\alpha} TU_{\alpha} \subseteq \bigcup_{A \in \cB} A \subseteq TM.
\end{align}
To show the second property of a basis, let $A \subseteq TU$ be open and $B \subseteq TV$ be open. Then $TU \cap TV = \pi^{-1}(U) \cap \pi^{-1}(V) = \pi^{-1}(U \cap V) = T(U \cap V)$ and so $A \cap B \subseteq T(U \cap V)$. Since $A \cap T(U \cap V)$ and $B \cap T(U \cap V)$ are both open in $T(U \cap V)$, $A \cap B = A \cap B \cap T(U \cap V)$ is open in $T(U \cap V)$.

\begin{lemma}
    $TM$ is a second countable Hausdorff space.
\end{lemma}
\begin{proof}
    Let $v,w \in TM$; so that $v \in T_{p}M$ and $w \in T_{q}M$, or $\pi(v) = p$ and $\pi(w) = q$. Two cases arise. If $p \neq q$, then since $M$ is Hausdorff, there exists open neighbourhoods $p \in U$ and $q \in V$ such that $U \cap V = \emptyset$. This simply implies $v \in \pi^{-1}(U) = TU$ and $w \in \pi^{-1}(V) = TV$, and $TU \cap TV = \emptyset$. For the case when $p = q$, both $v,w \in T_{p}M$. Since this $T_{p}M$ identifies with $\R^{n}$, $n$ being the dimension of $M$, we can pick open neighbourhoods $U_{1},U_{2} \subseteq T_{p}M$ of $v$ and $w$, respectively, such that $U_{1} \cap U_{2} = \emptyset$. If we now pick a coordinate chart $W$ of $p \in M$, then $W \times U_{1}$ and $W \times U_{2}$ are two open sets in $TM$ that separate $v$ and $w$, with $W \times U_{1} \cap W \times U_{2} = \emptyset$.

    Let $\{B_{i}\}_{i \geq 1}$ be a countable basis of $M$, and let $\{(U_{\alpha},\varphi_{\alpha})\}_{\alpha \in \Lambda}$ be an atlas. For $p \in U_{\alpha}$, there exists $p \in B_{p,\alpha} \subseteq U_{\alpha}$. Varying $p$ and $\alpha$, we get $\{B_{p,\alpha}\}_{p \in M, \, \alpha \in \Lambda} \subseteq \{B_{i}\}_{i \geq 1}$. So, there exists a countable basis of $M$ consisting of coordinate charts for $M$. Since $TU_{i}$ identifies with $\varphi_{i}(U_{i}) \times \R^{n}$ for any $i$, it is second countable and has a countable basis $\cB_{i}$. Therefore, $\cB = \bigcup_{i=1}^{\infty} \cB_{i}$ is a countable basis of $TM$.
\end{proof}

Finally, we show that $\{(TU_{\alpha},\tilde{\varphi}_{\alpha})\}$ gives an atlas consisting of smooth compatible charts on $TM$. Let $(U,\varphi) = (U,x^{1}\ldots,x^{n})$ and $(V,\psi) = (V,y^{1},\ldots,y^{n})$ be two charts in $M$ such that $U \cap V \neq \emptyset$. Let $p \in U \cap V$ with $v \in T_{p}M$. $\{\frac{\partial}{\partial x^{1}}\big|_{p},\ldots,\frac{\partial}{\partial x^{n}}\big|_{p}\}$ and $\{\frac{\partial}{\partial y^{1}}\big|_{p},\ldots,\frac{\partial}{\partial y^{n}}\big|_{p}\}$ are two bases for $T_{p}M$, so $v = \sum_{i=1}^{n} a^{i} \frac{\partial}{\partial x^{i}}\big|_{p} = \sum_{i=1}^{n} b^{i} \frac{\partial}{\partial y^{i}}\big|_{p}$. Applying $y^{i}$ on both sides gives
\begin{align}
    b^{i} = \sum_{j=1}^{n} \frac{\partial y^{i}}{\partial x^{j}}(p) a^{j}.
\end{align}
This is a smooth transformation, showing that switching from one set of coordinates to another is smooth. Moreover, the $TU_{\alpha}$'s cover $TM$. This is a valid atlas.

As a consequence of this, the projection map $\pi:TM \to M$ is a smooth map, since $\varphi(U) \times \R^{n} \cong TU \xrightarrow{\pi} U \cong \varphi(U)$ takes the first $n$ coordinates of $\varphi(U) \times \R^{n}$, which is a smooth map.

\begin{definition}
    A \eax{vector field} $X$ on a smooth manifold $M$ is a map $X:M \to TM$ such that $X(p) \in T_{p}M$ for all $p \in M$. In other words, a vector field is a smooth assignment of a tangent vector to each point in the manifold. It is termed smooth if the map $X$ is smooth, meaning that for any chart $(U,\varphi)$ on $M$, the composition $\tilde{\varphi} \circ X|_{U}: U \to \varphi(U) \times \R^{n}$ is smooth.
\end{definition}
Such a vector field $X$ is a section of the tangent bundle $TM$, as it assigns to each point $p \in M$ a vector in the corresponding tangent space $T_{p}M$. The projection map $\pi:TM \to M$ turns out to be the left inverse of $X$, since $\pi(X(p)) = p$ for all $p \in M$.

\begin{definition}
    Let $M$ be a smooth manifold of dimension $n$. A \eax{smooth vector bundle} on $M$ of rank $r$ is a triple $(E,M,\pi)$, where $E$ is a smooth manifold of dimension $n+r$, and $\pi:E \to M$ is a smooth surjective map such that for each $p \in M$, the fiber $\pi^{-1}(p)$ is an $r$-dimensional real vector space. Additionally, $\pi$ satisfies the local triviality condition: for every $p \in M$, there exists an open neighborhood $U$ of $p$ in $M$ and a diffeomorphism $\phi_{U}: E_{U} \defeq \pi^{-1}(U) \to U \times \R^{r}$ such that
    \begin{align}
        \pi|_{E_{U}} = \text{pr}_{1} \circ \phi_{U},
    \end{align}
    where $\text{pr}_{1}: U \times \R^{r} \to U$ is the projection onto the first factor, and $\varphi_{U}|_{\pi^{-1}(p)}: \pi^{-1}(p) \to \{p\} \times \R^{r}$ is a linear isomorphism of real vector spaces for each $p \in U$. Here, $E$ is termed the \eax{total space} and $M$ the \eax{base space} of the vector bundle.
\end{definition}

\begin{definition}
    A \eax{morphism} between two smooth vector bundles $(E,M,\pi)$ and $(F,N,\varphi)$ is a pair of smooth maps $(\tilde{f},f)$, where $\tilde{f}:E \to F$ and $f:M \to N$, such that $\varphi \tilde{f} = f \pi$, and for each $p \in M$, the map $\tilde{f}|_{\pi^{-1}(p)}: \pi^{-1}(p) \to \varphi^{-1}(f(p))$ is a linear map of real vector spaces. 
\end{definition}
As a special case, if $M = N$, and $f = \text{id}_{M}$, then $\tilde{f}$ is called a \eax{bundle map} over $M$. If $\tilde{f}$ is a diffeomorphism, then $(\tilde{f},\text{id}_{M})$ is an isomorphism of vector bundles, and the two bundles are said to be isomorphic. In particular, if $(E,M,\pi)$ is isomorphic to $(M \times \R^{r},M,\text{pr}_{1})$, then $(E,M,\pi)$ is called a \eax{trivial vector bundle} of rank $r$.
\subsection{Sections and Frames}
\textit{September 2nd.}

As earlier with a vector field, we have the following.

\begin{definition}
    A map $s:M \to E$ is called a \eax{section} of the vector bundle $(E,M,\pi)$ if $\pi(s(p)) = p$ for all $p \in M$. In other words, a section assigns to each point in the base space a vector in the corresponding fiber. A section is termed smooth if the map $s$ is smooth, meaning that for any chart $(U,\varphi)$ on $M$, the composition $\phi_{U} \circ s|_{U}: U \to U \times \R^{r}$ is smooth.
\end{definition}

We use $\Gamma(E) = \Gamma(M,E)$ to denote the set of all (global) smooth sections of the vector bundle $(E,M,\pi)$. For any open subset $U \subseteq M$, we denote by $\Gamma(U,E)$ the set of all smooth sections of $(E,M,\pi)$ over $U$; that is, smooth maps $s:U \to E$ such that $\pi(s(p)) = p$ for all $p \in U$. Moreover, for $s,t \in \Gamma(U,E)$, we define the sum $s+t \in \Gamma(U,E)$ by $(s+t)(p) = s(p) + t(p)$ for all $p \in U$; this sum is indeed in $\Gamma(U,E)$ since $\pi((s+t)(p)) = \pi(s(p) + t(p)) = \pi(s(p)) = p$, and $s+t$ can also be shown to be smooth. For $f \in C^{\infty}(U)$, we define the product $fs \in \Gamma(U,E)$ by $(fs)(p) = f(p)s(p)$ for all $p \in U$; this product is also in $\Gamma(U,E)$ since $\pi((fs)(p)) = \pi(f(p)s(p)) = \pi(s(p)) = p$, and $fs$ can also be shown to be smooth. Therefore, $\Gamma(U,E)$ is a module over the ring $C^{\infty}(U)$, and $\Gamma(E)$ is a module over the ring $C^{\infty}(M)$.

\begin{definition}
    A \eax{frame} of a vector bundle $(E,M,\pi)$ over an open subset $U \subseteq M$ is a collection of $r$ sections $s_{1},\ldots,s_{r} \in \Gamma(U,E)$ such that for each $p \in U$, the set $\{s_{1}(p),\ldots,s_{r}(p)\}$ forms a basis of the fiber $E_{p} = \pi^{-1}(p)$. The frame is termed smooth if each section $s_{i}$ is smooth. 
\end{definition}

Consider the tangent bundle $(TM,M,\pi)$ of a smooth manifold $M$ of dimension $n$. Let $(U,\varphi) = (U,x^{1},\ldots,x^{n})$ be a chart on $M$. Then the collection of sections $s_{i}:U \to TM$ defined by $s_{i}(p) = \frac{\partial}{\partial x^{i}}\big|_{p}$ for $i = 1,\ldots,n$ forms a smooth frame of the tangent bundle over $U$.

\begin{lemma}
    Let $(E,M,\pi)$ be a vector bundle, and let $U \subseteq M$ be an open subset. Suppose $s_{1},\ldots,s_{r}$ is a smooth frame of $E$ over $U$, and $s = \sum_{i=1}^{r} c^{i} s_{i}$. Then $s$ is a smooth section of $E$ over $U$ if and only if each coefficient function $c^{i}:U \to \R$ is smooth.
\end{lemma}
\begin{proof}
    The ``if'' direction is clear from the definition of smoothness. For the ``only if'' direction, let $s$ be a smooth section of $E$ over $U$ such that $s = \sum_{i=1}^{r} c^{i} s_{i}$. Let $p \in U$. We show that each $c^{i}$ is smooth at $p$. Since $s_{1},\ldots,s_{r}$ is a smooth frame of $E$ over $U$, there exists an open neighborhood $V \subseteq U$ of $p$ and a diffeomorphism $\phi_{V}: E_{V} \to V \times \R^{r}$ such that $\pi|_{E_{V}} = \text{pr}_{1} \circ \phi_{V}$, and $\phi_{V}|_{\pi^{-1}(q)}: \pi^{-1}(q) \to \{q\} \times \R^{r}$ is a linear isomorphism for each $q \in V$. Let $\phi_{V}(s(q)) = (q, (c^{1}(q),\ldots,c^{r}(q)))$ for all $q \in V$. Since $\phi_{V} \circ s|_{V}: V \to V \times \R^{r}$ is smooth, it follows that each component function $c^{i}: V \to \R$ is smooth. Therefore, each coefficient function $c^{i}$ is smooth at $p$, and since $p$ was arbitrary, each $c^{i}$ is smooth on all of $U$.
\end{proof}

\textit{September 4th.}
\begin{definition}
    Let $f:N \to M$ be a smooth map between smooth manifolds, such that $f$ is also injective and an immersion. Then $f(N)$, equipped with the topology borrowed from $N$, is termed an \eax{immersed submanifold}.
\end{definition}
In such a case, $f(N)$ is diffeomorphic to $N$. Note, however, that $f(N)$ need not be a manifold in the subspace topology of $M$; the subspace topology on $f(N)$ is weaker than the topology on $f(N)$ borrowed from $N$.

\begin{definition}
    A smooth map $f:N \to M$ between smooth manifolds is called an \eax{embedding} if $f$ is injective and an immersion, and the subspace topology from $M$ and the topology inherited from $N$ are the same on $f(N)$.
\end{definition}

\begin{theorem}
    If $S \subseteq M$ is a regular submanifold of $M$, then the inclusion map $i:S \to M$ is an embedding. Conversely, if $f:N \to M$ is an embedding, then $f(N)$ is a regular submanifold of $M$.
\end{theorem}
\begin{proof}
    Note that $i:S \to M$ is a smooth map, taking $(x^{1},x^{2},\ldots,x^{s}) \mapsto (x^{1},x^{2},\ldots,x^{s},0,\ldots,0)$, is injective and an immersion. Also, the topologies trivially coincide. Hence, $i$ is an embedding. For the converse, suppose $f:N \to M$ is an embedding. Let $p \in N$, with $f(p) \in f(N)$. Since $f$ is an immersion (at $p$), there exist charts $(U,\varphi) = (U,x^{1},\ldots,x^{n})$, centred at $p$, and $(V,\psi) = (V,y^{1},\ldots,y^{m})$, centred at $f(p)$, such that $f(U) \subseteq V$, and $f$ acts as $(x^{1},\ldots,x^{n}) \to (x^{1},\ldots,x^{n},0,\ldots,0)$. Now, since $f(U) \subseteq f(N) \cap V$, we will change $V$ to be the smallest open set $W \subseteq V \subseteq M$ such that $f(U) = f(N) \cap W$; since $U$ is open in $N$ if and only if $f(U)$ is open in $f(N)$ if and only if $f(U) = V' \cap f(N)$ where $V'$ is open in $M$, we can set $W = V \cap V' \subseteq V$, so that $f(U) = W \cap f(N)$.
\end{proof}

\begin{theorem}
    Let $f:N^{n} \to M^{m}$ be a smooth map between smooth manifolds. Suppose $S^{s} \subseteq M$ is a regular submanifold of $M$, and $f(N) \subseteq S$; this induces a set-theoretic factoring map $\tilde{f}:N \to S$ such that $i \circ \tilde{f} = f$. Then $\tilde{f}$ is a smooth map between smooth manifolds.
\end{theorem}
\begin{proof}
    With $p \in N$, and $f(p) \in S \subseteq M$, there exists a chart $(V,\psi) = (V,y^{1},\ldots,y^{m})$, centred at $f(p)$ in $M$, such that $S \cap V = \{(y^{1},\ldots,y^{m}) \in V \mid y^{s+1}=\cdots=y^{m}=0\}$. There also exits a chart $(U,\varphi) = (U,x^{1},\ldots,x^{n})$, centred at $p \in N$, such that $f(U) \subseteq V$ and $f$ is given by $(x^{1},\ldots,x^{n}) \mapsto (f^{1},f^{2},\ldots,f^{m})$, where each $f^{i}$ is smooth on $U$. But this means $f^{s+1}(x^{1},\ldots,x^{n}) = \cdots = f^{m}(x^{1},\ldots,x^{n}) = 0$ for all $(x^{1},\ldots,x^{n}) \subseteq U$. In $S$, we have a smooth chart given by $(V \cap S, \varphi_{S})$ via which, the map $\tilde{f}$ on $U$ is given as $(x^{1},\ldots,x^{n}) \mapsto (f^{1},\ldots,f^{s})$; $\tilde{f}$ is smooth.
\end{proof}

\section{Bump Functions and Partitions of Unity}
\textit{September 18th.}

\begin{definition}
    A \eax{bump function} is a smooth non-negative function $f$ on a smooth manifold $M$ such that $f$ has compact support, that is, the closure of the set $\{q \in M \mid f(q) \neq 0\}$ is compact in $M$. A bump function is said to be adapted to an open set $U \subseteq M$ if the support of $f$ is contained in $U$.
\end{definition}

If $p \in U \subseteq M$ is an open set, then there exists a bump function $f$ adapted to $U$ such that $f(x) = 1$ in a neighbourhood of $p$ contained in $U$.

\begin{definition}
    A collection of subsets $\{A_{\alpha}\}_{\alpha \in \Lambda}$ of a topological space $S$ is said to be \eax{locally finite} if for all $p \in S$, there exists a neighbourhood $V$ of $P$ such that $\{\alpha \in \Lambda \mid A_{\alpha} \cap V \neq \emptyset\}$ is finite. In other words, each point in $S$ has a neighbourhood that intersects only finitely many of the sets $A_{\alpha}$.
\end{definition}

\begin{definition}
    A (smooth) \eax{partition of unity} is a collection of smooth functions $\{\rho_{\alpha}:M \to \R\}_{\alpha \in \Lambda}$ on a smooth manifold $M$ such that the following conditions hold:
    \begin{enumerate}
        \item $\{\supp{\rho_{\alpha}}\}_{\alpha \in \Lambda}$ is a locally finite collection of subsets of $M$.
        \item $\sum_{\alpha \in \Lambda} \rho_{\alpha}(p) = 1$ for all $p \in M$.
    \end{enumerate}
\end{definition}
Note that the sum in the second condition is well-defined since for each $p \in M$, there exists a neighbourhood $V$ of $p$ such that $\{\alpha \in \Lambda \mid \supp{\rho_{\alpha}} \cap V \neq \emptyset\}$ is finite, and so only finitely many terms in the sum are non-zero.

\begin{theorem}
    Let $M$ be a smooth manifold, and let $\{U_{\alpha}\}_{\alpha \in \Lambda}$ be an open cover of $M$. Then the following hold true.
    \begin{enumerate}
        \item There exists a countable smooth partition of unity $\{\varphi_{i}\}_{i=1}^{\infty}$ on $M$ such that for each $i$, $\supp(\varphi_{i})$ is compact and contained in some $U_{\alpha}$.
        \item There exists a smooth partition of unity $\{\rho_{\alpha}\}_{\alpha \in \Lambda}$ on $M$ such that for each $\alpha$, $\supp(\rho_{\alpha})$ is contained in $U_{\alpha}$. 
    \end{enumerate}
\end{theorem}
In the second statement, we say that the partition of unity is \eax{subordinate} to the open cover $\{U_{\alpha}\}_{\alpha \in \Lambda}$. Before we prove the theorem, we require a few results from topology.

Firstly, since $M$ is a smooth manifold, it is second countable and has a countable basis $\{B_{i}\}_{i \geq 1}$. Define $\cS = \{B_{i} \mid \overline{B_{i}} \text{ is compact}\}$. Then $\cS$ is a countable basis of $M$ consisting of open sets with compact closure. We shall show that $\cS$ also happens to be a basis of $M$; to this end, let $p \in U \subseteq M$ be a point in an open set $U$. Moreover, let $V$ be an open neighbourhood of $p$ such that $\overline{V}$ is compact and $\overline{V} \subseteq U$. Since $\cB$ is a basis of $M$, there exists $B_{i} \in \cB$ such that $p \in B_{i} \subseteq V$. But then $\overline{B_{i}} \subseteq \overline{V}$ is compact, and so $B_{i} \in \cS$. Therefore, $\cS$ is a basis of $M$ consisting of open sets with compact closure. Thus, any smooth manifold has a countable basis of open sets with compact closure.

Secondly, given any smooth manifold $M$, there exists a sequence of open sets $\{V_{i}\}_{i \geq 1}$ such that $\overline{V_{i}}$ is compact and $\overline{V_{i}} \subseteq V_{i+1}$ for all $i$, and $M = \bigcup_{i=1}^{\infty} V_{i}$. Indeed, let $\{B_{i}\}_{i=1}^{\infty}$ be a countable basis of $M$ consisting of open sets with compact closure. Define $V_{1} = B_{1}$, so that $V_{1} \subseteq \overline{V}_{1}$ is compact. Thus, it must be contained in a finite union of the open cover $\{U_{\alpha}\}_{\alpha \in \Lambda}$; that is, there exists a finite set $A_{1} \subseteq \Lambda$ such that $\overline{V}_{1} \subseteq \bigcup_{\alpha \in A_{1}} U_{\alpha}$; for simplicity, suppose $\overline{V}_{1} \subseteq B_{1} \cup B_{2} \cup \cdots \cup B_{i_{1}}$ for some $i_{1} > 1$ such that $i_{1}$ is the smallest integer with this property. Now, let $V_{2} = B_{1} \cup B_{2} \cup \cdots \cup B_{i_{1}}$, an open set with compact closure. Thus, $V_{1} \subseteq \overline{V}_{1} \subseteq V_{2}$, and $\overline{V}_{2}$ is compact. Repeating the same argument, we have $\overline{V}_{2} \subseteq B_{1} \cup B_{2} \cup \cdots \cup B_{i_{1}} \cup \cdots \cup B_{i_{2}}$ for some $i_{2} > i_{1}$ such that $i_{2}$ is the smallest integer with this property. Let $V_{3} = B_{1} \cup B_{2} \cup \cdots \cup B_{i_{2}}$, an open set with compact closure. Thus, $V_{2} \subseteq \overline{V}_{2} \subseteq V_{3}$, and $\overline{V}_{3}$ is compact. Continuing in this manner, we get a sequence of open sets $\{V_{i}\}_{i=1}^{\infty}$ such that $\overline{V}_{i}$ is compact and $\overline{V}_{i} \subseteq V_{i+1}$ for all $i$, and $M = \bigcup_{i=1}^{\infty} V_{i}$. If we define $V_{-1} = V_{0} = \emptyset$, we obtain a sequence as
\begin{align}
    V_{-1} \subseteq \overline{V}_{-1} \subseteq V_{0} \subseteq \overline{V}_{0} \subseteq V_{1} \subseteq \overline{V}_{1} \subseteq V_{2} \subseteq \overline{V}_{2} \subseteq V_{3} \subseteq \cdots
\end{align}
From here, we can infer that
\begin{align}
    \overline{V}_{i+1} \setminus V_{i} \subseteq V_{i+2} \setminus \overline{V}_{i-1}
\end{align}
where $\overline{V}_{i+1} \setminus V_{i}$ is compact and $\overline{V}_{i-1} \subseteq V_{i+2}$ is open. We shall use ths construction in the proof of the theorem.

\begin{proof}
    Start with the construction of the sequence of open sets $\{V_{i}\}_{i=-1}^{\infty}$ as above, where $V_{-1} = V_{0} = \emptyset$, and $\overline{V}_{i} \subseteq V_{i+1}$ for all $i$. Let $p \in U_{\alpha}$ for some $\alpha \in \Lambda$. Then there exists $i$ such that $p \in U_{\alpha} \cap (V_{i+2}-\overline{V}_{i-1})$. Since $U_{\alpha} \cap (V_{i+2}-\overline{V}_{i-1})$ is open, there exists a smooth bump function $\rho_{p}:M \to \R$ such that $\supp(\rho_{p}) \subseteq U_{\alpha} \cap (V_{i+2}-\overline{V}_{i-1})$ and $\rho_{p} > 0$ in some neighbourhood $W_{p}$ of $p$. The collection $\{W_{p}\}_{p \in \overline{V}_{i+1} \setminus V_{i}}$ is then an open cover of the compact set $\overline{V}_{i+1} \setminus V_{i}$; refine this cover to a finite subcover consisting of $W_{1},\ldots,W_{m}$, and let $\rho_{1},\ldots,\rho_{m}$ be the corresponding bump functions. Repeating this for every $i$, we obtain, for each $\overline{V}_{i+1}\setminus V_{i}$, a finite cover $W_{1}^{i},W_{2}^{i},\ldots,W_{m(i)}^{i}$ of $\overline{V}_{i+1}\setminus V_{i}$, and corresponding bump functions $\rho_{1}^{i},\rho_{2}^{i},\ldots,\rho_{m(i)}^{i}$ such that $\supp(\rho_{j}^{i}) \subseteq U_{\alpha} \cap (V_{i+2}-\overline{V}_{i-1})$ for some $\alpha \in \Lambda$, and $\rho_{j}^{i} \equiv 1$ in $W_{j}^{i}$ for all $j = 1,\ldots,m(i)$. Now, let $\{\rho_{k}\}_{k=1}^{\infty}$ be the countable collection of all these bump functions, and $\{W_{j}^{i}\}$ be the corresponding open sets. Note that the following properties hold:
    \begin{enumerate}
        \item $\supp(\rho_{j}^{i})$ is compact for all $i,j$.
        \item $\supp(\rho_{j}^{i}) \subseteq U_{\alpha_{ij}} \cap (V_{i+2}-\overline{V}_{i-1})$ for some $\alpha_{ij} \in \Lambda$.
        \item $\overline{V}_{i+1} \setminus V_{i} \subseteq \bigcup_{j=1}^{m(i)} W_{j}^{i}$ for all $i$.
        \item $\rho_{j}^{i} > 0$ on $W_{j}^{i}$ for all $i,j$.
    \end{enumerate}
    Since $\supp(\rho_{j}^{i}) \subseteq V_{i+2} \setminus \overline{V}_{i-1}$, we must have $\supp(\rho_{j}^{i}) \cap \overline{V}_{i-1} = \emptyset$ and $\supp(\rho_{j}^{i}) \cap (M \setminus V_{i+2}) = \emptyset$. Therefore, for each $p \in M$, there exists $i$ such that $p \in V_{i+2} \setminus \overline{V}_{i-1}$, and so $p$ can only belong to the support of bump functions $\rho_{j}^{k}$ for $k = i-1,i,i+1$. This shows that the collection $\{\supp(\rho_{j}^{i})\}$ is locally finite. Finally, we define
    \begin{align}
        \varphi_{k} = \frac{\rho_{k}}{\sum_{j=1}^{\infty} \rho_{j}}
    \end{align}
    where the denominator is smooth and positive, and only finitely many terms are non-zero in a neighbourhood of any point. Then $\{\varphi_{k}\}_{k=1}^{\infty}$ is a countable smooth partition of unity on $M$ such that for each $k$, $\supp(\varphi_{k})$ is compact and contained in some $U_{\alpha}$. This proves the first statement. For the second statement, we can define $\rho_{\alpha} = \sum_{k \in I_{\alpha}} \varphi_{k}$, where $I_{\alpha} = \{k \mid \supp(\varphi_{k}) \subseteq U_{\alpha}\}$. Then $\{\rho_{\alpha}\}_{\alpha \in \Lambda}$ is a smooth partition of unity on $M$ subordinate to the open cover $\{U_{\alpha}\}_{\alpha \in \Lambda}$.
\end{proof}

\begin{lemma}
    Let $M$ be a smooth manifold, and $X$ be a vector field on $M$. Then $X$ is smooth on $M$ if and only if for any chart $(U,\varphi)$ on $M$, with $X = \sum_{i=1}^{n} a^{i} \frac{\partial}{\partial x^{i}}$ on $U$, each coefficient function $a^{i}:U \to \R$ is smooth.
\end{lemma}
\begin{proof}
    Suppose $X$ was smooth on $M$, or $X:M \to TM$ is a smooth map. Then for any chart $(U,\varphi)$ on $M$, the composition $\tilde{\varphi} \circ X|_{U}: U \to \varphi(U) \times \R^{n}$ is smooth, and so each component function $a^{i}:U \to \R$ is smooth. Conversely, suppose that for any chart $(U,\varphi)$ on $M$, with $X = \sum_{i=1}^{n} a^{i} \frac{\partial}{\partial x^{i}}$ on $U$, each coefficient function $a^{i}:U \to \R$ is smooth. Then for any chart $(U,\varphi)$ on $M$, the composition $\tilde{\varphi} \circ X|_{U}: U \to \varphi(U) \times \R^{n}$ is smooth, and so $X:M \to TM$ is a smooth map.
\end{proof}

\begin{proposition}
    Let $X$ be a vector field on a smooth manifold $M$. Then $X$ is smooth if and only if for any smooth function $f \in C^{\infty}(M)$, the function $X(f):M \to \R$ defined by $X(f)(p) = X_{p}(f)$ is smooth.
\end{proposition}
\begin{proof}
    Suppose $X$ is smooth. Then for any smooth function $f \in C^{\infty}(M)$, the function $X(f):M \to \R$ defined by $X(f)(p) = X_{p}(f)$ is smooth. Conversely, suppose that for any smooth function $f \in C^{\infty}(M)$, the function $X(f):M \to \R$ defined by $X(f)(p) = X_{p}(f)$ is smooth. Let $(U,\varphi)$ be a chart on $M$, with $\varphi = (x^{1},\ldots,x^{n})$. Then on $U$, we have $X = \sum_{i=1}^{n} a^{i} \frac{\partial}{\partial x^{i}}$, where each coefficient function $a^{i}:U \to \R$ is given by $a^{i}(p) = X_{p}(x^{i})$. Since each coordinate function $x^{i}$ is smooth, it follows that each coefficient function $a^{i}:U \to \R$ is smooth. Therefore, by the previous lemma, $X$ is smooth on $M$.
\end{proof}
Moreover, since $X_{p}$ is a $\R$-linear derivation from $C^{\infty}(M)$ to $\R$, we get that $X$ is a derivation on $C^{\infty}(M)$; that is, for any $f,g \in C^{\infty}(M)$, we have
\begin{align}
    X(fg)(p) = X_{p}(fg) = f(p)X_{p}(g) + g(p)X_{p}(f) = (fX(g) + gX(f))(p) \implies X(fg) = fX(g) + gX(f).
\end{align}


\section{Integral Curves}
The following theorem will be useful in the study of vector fields and their flows on smooth manifolds.
\begin{theorem}
    Let $p_{0} \in V \subseteq \R^{n}$ be an open set, and $f:V \to \R^{n}$ be a smooth map. Then the differential equation $\frac{dy}{dt} = f(y)$ with initial condition $y(0) = p_{0}$ has a unique smooth solution $y:(a(p_{0}),b(p_{0})) \to V$ where $(a(p_{0}),b(p_{0}))$ is the maximal open interval containing $0$ such that $y$ is defined on $(a(p_{0}),b(p_{0}))$.
\end{theorem}

\begin{definition}
    Let $X$ be a smooth vector field on a smooth manifold $M$, and let $p \in M$. An \eax{integral curve} of $X$ is a smooth curve $c:(a,b) \to M$ such that $c'(t) = X_{c(t)}$ for all $t \in (a,b)$. We usually presume $0 \in (a,b)$; if $c(0) = p$, we say that $c$ is an integral curve of $X$ starting at $p$.

    A \eax{maximal integral curve} is an integral curve whose domain cannot be extended to a larger interval while still satisfying the integral curve condition. 
\end{definition}

We look at a few examples.
\begin{example}
    Consider the smooth manifold $M = \R^{2}$ and the smooth vector field $X = x \frac{\partial}{\partial x} + y \frac{\partial}{\partial y}$. Let $p_{0} = (1,0) \in \R^{2}$. Then, an integral curve $c:(a,b) \to \R^{2}$ of $X$ starting at $p_{0}$ satisfies the differential equation $c'(t) = (\dot{x}(t),\dot{y}(t)) = (-y(t),x(t)) = X_{c(t)}$ with initial condition $c(0) = (1,0)$. Solving this system of differential equations, we get
    \begin{align}
        c(t) = (\cos t, \sin t).
    \end{align}
    This integral curve is a circle centered at the origin with radius 1, and it is defined for all $t \in \R$. If, instead, we had chosen $p_{0} = (x_{0},y_{0}) \in \R^{2}$, then the integral curve of $X$ starting at $p_{0}$ would be given by
    \begin{align}
        c(t) = (x_{0}\cos t - y_{0}\sin t, x_{0}\sin t + y_{0}\cos t),
    \end{align}
    which is a circle centered at the origin with radius $\sqrt{x_{0}^{2} + y_{0}^{2}}$, and it is also defined for all $t \in \R$. Moreover, one can verify that $c_{s}(c_{t}) = c_{s+t}$ for all $s,t \in \R$, where $c_{t}(p_{0})$ denotes the integral curve of $X$ starting at $p_{0}$ evaluated at time $t$. $c_{t}$ is a diffeomorphism of $\R^{2}$ for each $t \in \R$, and the map
    \begin{align}
        \R \to \mathrm{Diff}(\R^{2}), \quad t \mapsto c_{t}
    \end{align}
    is a smooth group homomorphism from $(\R,+)$ to the group of diffeomorphisms of $\R^{2}$, where the group operation on $\mathrm{Diff}(\R^{2})$ is composition of maps. 
\end{example}

From the above, any group homomorphism from $(\R,+)$ to the group of diffeomorphisms of a smooth manifold $M$ is called a one-parameter group of diffeomorphisms of $M$.

\begin{corollary}
    Let $X$ be a smooth vector field on a smooth manifold $M$, with $p \in M$. Then there exists an integral curve $c:(a,b) \to M$ of $X$ such that $0 \in (a,b)$ with $c(0) = p$, and this integral curve is unique.
\end{corollary}

\begin{theorem}
    Let $V \subseteq \R^{n}$ be open, and $f:V \to \R^{n}$ be a smooth map. Then for each point $p_{0} \in V$, there exists a neighbourhood $W$ of $p_{0}$ in $V$, an $\epsilon > 0$, and a smooth function $y:(-\epsilon,\epsilon) \times W \to V$ such that
    \begin{align}
        \frac{\partial y}{\partial t}(t,q) = f(y(t,q)), \qquad y(0,q) = q
    \end{align}
    for all $(t,q) \in (-\epsilon,\epsilon) \times W$.
\end{theorem}
\begin{corollary}
    Let $X$ be a smooth vector field on a smooth manifold $M$, and let $U$ be a chart around $p \in M$. Then there exists a neighbourhood $W$ of $p$ in $U$, an $\epsilon > 0$, and a smooth map $F:(-\epsilon,\epsilon) \times W \to U$ such that for all $q \in W$, $F(t,q)$ is an integral curve of $X$ with initial point $q$ ($F(0,q) = q)$.
\end{corollary}
In the above, we use the notation $F(t,q) = F_{t}(q)$. Suppose $s,t \in (-\epsilon,\epsilon)$, and $q \in W$, and that $F_{t}(F_{s}(q))$ and $F_{t+s}(q)$ are both defined. Then we have $F_{t}(F_{s}(q)) = F_{t+s}(q)$. This is because both sides are integral curves of $X$ starting at $q$ at time $0$, and by uniqueness of integral curves, they must coincide. $F$, above, is termed as a \eax{local flow} generated by the vector field $X$. $F_{t}(q)$ is then a flow line of $q$ of the local flow $F$. If $F$ can be defined as $F:\R \times M \to M$, then $F$ is called a \eax{global flow} generated by the vector field $X$. $X$ is called a \eax{complete vector field} if its flow is global.

\begin{definition}
    In general, a \eax{local flow} about a point $p \in U \subseteq M$ is a smooth map $F:(-\epsilon,\epsilon) \times W \to U$, where $W$ is a neighbourhood of $p$ in $U$ and $\epsilon > 0$, such that
    \begin{enumerate}
        \item $F_{0}(q) = q$ for all $q \in W$, and
        \item $F_{t}(F_{s}(q)) = F_{t+s}(q)$ whenever both sides are defined.
    \end{enumerate}
\end{definition}